% Source: https://artofproblemsolving.com/community/c3h3181940_2022_abmc_online_october_contest__actonboxborough_math_competition
% Stable identity: aops-topic-3181940-q12
The triangle area is 6. The hypotenuse is distance $\dfrac{12}{5}$ from $B$, while the radius is $8\dfrac{\sqrt3}{5}$. The circle cuts the hypotenuse in a chord subtending angle $2\arccos(\dfrac{\sqrt3}{2})=\dfrac{\pi}{3}$. The small disk segment outside the triangle has area $\dfrac{1}{2}r^2(\dfrac{\pi}{3})-\dfrac{1}{2}r^2\sin(\dfrac{\pi}{3})=\dfrac{32\pi-48\sqrt3}{25}$. Since the radius is less than both legs, the disk initially contributes one quarter disk inside the angle at $B$, of area $48\dfrac{\pi}{25}$. Therefore the area outside the disk is $6-48\dfrac{\pi}{25}+\dfrac{32\pi-48\sqrt3}{25}=\dfrac{150-48\sqrt3-16\pi}{25}$. Hence $\boxed{242}$.
\begin{asy}
usepackage("amsmath");
size(7cm);
pair B=(0,0),A=(3,0),C=(0,4);draw(A--B--C--cycle);draw(arc(B,8*sqrt(3)/5,0,90),blue);label("$A$",A,SE);label("$B$",B,SW);label("$C$",C,N);
\end{asy}
